\chapter{План работ}

\subsection*{Описание проекта:}
Цель проекта: Разработка ПО для минимизации высоты укладки 3D-грузов на паллету с помощью гибридного алгоритма (имитация отжига + жадная эвристика).

Представляет интерес решение нетривиальной задачи по упаковке и реализация алгоритма с его визуализацией.

Объект исследования: Процесс укладки разногабаритных коробок на стандартную грузовую паллету.

Предмет исследования: Алгоритмы решения задачи трехмерной упаковки (3D Bin Packing) и методы их оптимизации.

Предполагается, что комбинирование метода имитации отжига (для поиска порядка) и жадного алгоритма (для размещения) позволит уменьшить итоговую высоту укладки по сравнению с простыми эвристиками.

Задача относится к классу комбинаторной оптимизации (3D-BPP). Точные методы решения требуют огромных вычислительных мощностей. Современный подход заключается в использовании гибридных алгоритмов, балансирующих между скоростью работы и качеством укладки.
\begin{table}[h!]
\centering
\caption{План выполнения проекта}
\label{tab:project_plan}
\begin{tabular}{|p{3cm}|p{6cm}|p{4cm}|p{3cm}|}
\hline
\textbf{Этап проекта} & \textbf{Описание работ} & \textbf{Ожидаемые результаты} & \textbf{Сроки выполнения} \\
\hline
1. Аналитический этап & Изучение проблематики и подготовка обзора литературы. & Литературный обзор, обоснование выбора гибридного алгоритма. & Декабрь 2025 \\
\hline
2. Проектирование & Создание структуры исследовательской модели и ее загрузка в GitHub. & Структура исследовательской модели & Декабрь 2025 \\
\hline
3. Программная реализация простой модели & Реализация простой модели для укладки одной паллеты на эвристиках. & Рабочая программа простой модели & Декабрь 2025 \\
\hline
4. Статья для простого  случая & Подготовка статьи . & Статья & Январь 2026 \\
\hline
5. Модель на основе методов оптимизации & Реализация модели на основе методов оптимизации, подготовка презентации. & Готовая модель & Март 2026 \\
\hline
6. Модель для множества паллет & Реализация модели для множества паллет. & Готовая модель для множества паллет & Март 2026 \\
\hline
7. Заключение & В заключение должны быть сформированы критерии для оценки эффективности предложенных алгоритмов и подготовлена публикация на русском или английском языке. & Критерии оценки эффективности и публикация на русском & Апрель 2026 \\
\hline
\end{tabular}
\end{table}

\chapter{Критерии оценивания проекта}
\begin{table}[h!]
\centering
\caption{Критерии оценки эффективности алгоритмов}
\label{tab:criteria}
\begin{tabular}{|p{4cm}|p{10cm}|}
\hline
\textbf{Критерий} & \textbf{Описание} \\
\hline
Корректность теоретических решений & Правильность построенной модели и предложенных алгоритмов \\
\hline
Программная реализация & Эффективность и качество программной реализации \\
\hline
Анализ результатов экспериментов & Глубина и корректность анализа экспериментальных данных \\
\hline
\end{tabular}
\end{table}
